{\heiti 摘要：}
自主水下航行器（下称航行器）已能在强流海区执行长航程任务，其航路规划依赖岸基系统发布的海流图。受水声通信带宽与上浮代价的限制，航行器在航行中往往只有一张数天前的粗分辨率海流图（下称粗图）；绕越禁航区时，从左侧还是右侧通过是不可逆的选择，选错一侧要多花数小时。现有方法难以判断粗图能否支撑这一选择：加密求解网格只能减小求解误差，纠正不了粗图本身的误差；集合规划需要在大量随机流场上求解，且只统计改变选边的成员比例，不衡量选错的代价；自适应采样以降低全场的估计误差为目标，测量的未必是决定选边的海流。为此，本文提出合理反例检验：在与粗图相容的合理流场集合中，搜索使另一侧更快的流场（合理反例），以坚持当前选择最多多花的航时（最大后悔值）判断粗图是否可信，并据此决定继续计算还是获取新信息。该方法有三个要点：一是用多年再分析数据中粗图的误差，标定合理流场集合的误差大小、空间相关长度与半径；二是利用包络定理沿最快航线累加航时对流场的导数，约71次航线求解即可完成搜索；三是两侧航时差小于求解误差时才加密网格，最大后悔值超过阈值时上浮获取新的海流图，不能上浮时改选最大后悔值较小的一侧。我们在黑潮、湾流、阿古拉斯回流区以及改带卫星测高地转流图的黑潮区共2532个测试任务上，将该方法与加密求解、集合规划、就地测流和总是上浮等做法比较。结果表明，计算量增加64倍几乎不降低选错率，换用当天0.5°海流图则使之降至7\%--10\%；按最大后悔值排在前20\%的任务集中了全部额外航时（比真实更快一侧多花的航时）的53\%--71\%，高于按改变选边的成员比例排序的36\%--52\%；在黑潮，梯度搜索以约二十九分之一的求解次数达到集合平均选边的效果；在湾流与阿古拉斯回流区，对判定不可信的任务上浮取图，使平均额外航时分别从3.56小时降至1.31小时、从6.87小时降至2.20小时。这表明，粗图缺失的信息无法靠端侧计算补足，合理反例检验为航行器在计算与观测之间的取舍提供了可计算的判据。

{\heiti 关键词：}海洋大数据；自主水下航行器；航路规划；不确定性量化；最大后悔值；合理反例检验

\vspace{1em}
\begin{center}
{\large\bfseries Plausible Counterexample Test: When an Autonomous Underwater Vehicle Should Compute and When It Should Observe Under a Coarse Current Map}
\end{center}
\vspace{0.5em}

{\bfseries Abstract:}
Autonomous underwater vehicles (AUVs) now carry out long missions in strong-current regions, and their route planning relies on current maps issued by shore-based systems. Because acoustic links are slow and surfacing interrupts the mission, a vehicle at sea usually holds only a coarse current map issued days earlier. When it passes a no-go zone, the choice between the left and the right side cannot be undone, and the wrong side costs hours. Existing methods cannot tell whether the map supports this choice: a finer solver reduces numerical error but not the error of the map itself; ensemble planning needs many sampled currents and counts how many members change the side rather than what a wrong side costs; adaptive sampling reduces the error of the whole field instead of measuring the currents that decide the side. We propose a plausible counterexample test. It searches the set of currents consistent with the map for one that makes the other side faster (a plausible counterexample), uses the largest extra time of keeping the current choice (the maximum regret) to judge whether the map can be trusted, and on that basis decides whether to compute more or to acquire new information. The test has three parts. The set of plausible currents is calibrated from the errors of the map in multi-year reanalysis data. The derivative of travel time with respect to the current is accumulated along the fastest route by the envelope theorem, so the search needs about 71 route solves. A finer grid is used only when the two sides differ by less than the numerical error, and when the maximum regret exceeds a threshold the vehicle surfaces for a new map or, if it cannot surface, takes the side with the smaller maximum regret. On 2532 test missions in the Kuroshio, the Gulf Stream, the Agulhas retroflection and the Kuroshio with a satellite altimetry map, we compare the test with finer solving, ensemble planning, in-situ current measurement and always surfacing. Sixty-four times more computation barely lowers the share of wrong sides, whereas a same-day 0.5° map lowers it to 7\%--10\%. The 20\% of missions ranked riskiest by maximum regret hold 53\%--71\% of the extra time (the time beyond that of the truly faster side), against 36\%--52\% when ranked by the share of ensemble members that change the side. In the Kuroshio, the gradient search matches ensemble averaging with about one twenty-ninth of the route solves. In the Gulf Stream and the Agulhas, surfacing on missions judged untrustworthy cuts the mean extra time from 3.56 to 1.31 hours and from 6.87 to 2.20 hours. Information missing from the map cannot be recovered by onboard computation, and the test gives a computable criterion for trading computation against observation.

{\bfseries Keywords:} ocean big data; autonomous underwater vehicle; path planning; uncertainty quantification; maximum regret; plausible counterexample test
